% !TEX program = xelatex
\documentclass[12pt,a4paper]{article}

\usepackage[margin=2.5cm]{geometry}   % 页边距
\usepackage{booktabs}                 % 表格三条线 \toprule \midrule \bottomrule
\usepackage{tabularx}                 % 表格自动占满页宽
\usepackage{array}                    % 定宽列 p{}
\usepackage{hyperref}                 % 交叉引用可点击
\usepackage{setspace}
\usepackage[scheme=plain]{ctex}
\onehalfspacing                       % 1.5 倍行距

\begin{document}

\section*{A-III. From GFS, MapReduce and Bigtable to Current Cloud Frameworks}

\subsubsection*{Original designs}

GFS is a file system designed for storing large-scale data. 
In practice, storing massive files often presents challenges: 
a single machine may lack sufficient capacity, the hardware used is typically inexpensive, and disks are prone to failure. 
To address the issue of capacity limitations on individual machines, GFS divides large files into smaller chunks. 
These chunks are distributed across the system, allowing each one to be placed, replicated, or migrated independently. 
Since this fragmentation means that individual machines have no way of knowing where all the parts of a complete file are located, 
GFS employs a master node that acts as an index, tracking the locations of all chunks and guiding clients to the appropriate machines. 
Once a client learns the location of a specific chunk, it contacts the corresponding machine to retrieve it. 
To handle the risk of hardware failure, GFS incorporates a replication mechanism: each chunk is copied across multiple machines. 
If the master node detects that the number of replicas for a given chunk has fallen below a defined threshold—signaling that a machine has failed—it automatically creates new replicas on other machines, thereby ensuring data reliability.


\subsection*{Current successors}

\end{document}